% !TEX TS-program = pdflatex
% Lettre de motivation — Cloud & DevOps — AFD TECH part of Accenture — STYLE SOFRECOM/AZIT
\documentclass[a4paper,11pt]{article}
\usepackage[utf8]{inputenc}
\usepackage[T1]{fontenc}
\usepackage[scaled=0.94]{helvet}
\renewcommand{\familydefault}{\sfdefault}
\usepackage[left=2.5cm,right=2.5cm,top=2cm,bottom=2cm]{geometry}
\usepackage{xcolor}
\usepackage[hidelinks]{hyperref}
\definecolor{navy}{RGB}{23,54,93}
\pagestyle{empty}
\setlength{\parindent}{0pt}
\setlength{\parskip}{0pt}
\hypersetup{pdftitle={Lettre de motivation - Baha Eddine Belhaj Mouldi - AFD TECH Cloud DevOps},pdfauthor={Baha Eddine Belhaj Mouldi}}
\begin{document}

\begin{flushright}
{\bfseries Baha Eddine Belhaj Mouldi}\\
Tunis, Tunisie\\
+216 55 128 982\\
\href{mailto:bahaeddine.belhajmouldi@gmail.com}{bahaeddine.belhajmouldi@gmail.com}\\
\href{https://linkedin.com/in/bahamouldi}{linkedin.com/in/bahamouldi}\\
\href{https://github.com/bahamouldi}{github.com/bahamouldi}
\end{flushright}

\vspace{0.8cm}

\begin{flushleft}
{\bfseries Mme Nargesse BENAYADA}\\
AFD TECH part of Accenture — People \& Project Specialist

\vspace{0.8cm}

Tunis, le 01 septembre 2026

\vspace{0.6cm}

{\bfseries Objet : Candidature — Cloud \& DevOps Engineer (AWS/Azure/GCP, Kubernetes, ArgoCD, Terraform) — AFD TECH}
\end{flushleft}

\vspace{0.6cm}

Madame,

Votre campagne de recrutement pour les trois piliers \textit{Cloud \& DevOps, Middleware/BI \& Database, Infrastructure Core} chez AFD TECH part of Accenture a retenu toute mon attention. Peu d'entreprises au Maroc articulent aussi clairement l'ensemble de la chaîne DevOps — du provisionning Terraform à l'observabilité Dynatrace, en passant par ArgoCD, Vault et CSPM — avec une exigence opérationnelle qui correspond exactement à mon parcours en production sur Kubernetes.

Diplômé de l'École Nationale d'Ingénieurs de Tunis en juin 2026 (Mention Très Bien, 65e/600+ au concours national, PFE noté 86/100), j'ai conçu, déployé et opéré BeeWAF — un système en production sur un cluster Kubernetes de 5 nœuds (3 masters + 2 workers, Calico) — avec un pipeline CI/CD complet : Jenkins en huit étapes et ArgoCD en approche GitOps, manifests Kubernetes complets (Deployment, Service, Ingress, PVC, Secrets, ConfigMap), probes et autoscaling. Ce projet m'a appris à gérer des systèmes distribués à haute performance, avec une exigence zéro downtime (maxUnavailable:0) et une latence overhead inférieure à 20ms, supervisé par Prometheus/Grafana (26 panels) et ELK.

Au-delà du socle Kubernetes, j'ai mis en œuvre l'Infrastructure as Code avec Terraform (remote state) et Ansible, l'observabilité avec ELK et Dynatrace, ainsi que les bonnes pratiques DevSecOps (RBAC, IAM, Vault, WAF, virtual patching de 37 CVE). Mon expérience de pipeline de détection via data mining — clustering K-Means (silhouette 0.787) sur logs AWS avec explicabilité SHAP — m'a appris que la valeur d'un pipeline ne tient pas à sa complexité, mais à sa capacité à être monitoré, maintenu et sécurisé en production. C'est cette exigence que j'ai prolongée en automatisant les déploiements et le reporting avec Python/Bash.

Habitué à travailler en environnement international et à collaborer en transverse avec les équipes développement, je suis à l'aise aussi bien dans le provisionning d'infrastructure que dans l'optimisation des processus de livraison. Je souhaite mettre cette polyvalence Cloud, Middleware et Infra Core au service de vos équipes.

Je me tiens à votre entière disposition pour un entretien afin de vous exposer plus en détail ma motivation et mon expérience, et vous prie d'agréer, Madame, l'expression de mes salutations distinguées.

\vspace{0.8cm}

\begin{flushright}
{\bfseries Baha Eddine Belhaj Mouldi}
\end{flushright}

\end{document}
